\documentclass[11pt]{article}
\usepackage[a4paper,margin=26mm]{geometry}
\usepackage{amsmath,amssymb,amsthm,mathtools}
\usepackage{graphicx}
\usepackage{microtype}
\usepackage[hidelinks]{hyperref}
\usepackage{enumitem}

\newtheorem{theorem}{Theorem}
\newcommand{\Ep}{\mathcal E_p}
\newcommand{\R}{\mathbb R}

\title{An anisotropic obstruction to every $q>p$\\
reverse affine Zhang improvement}
\author{Candidate counterexample for arXiv:2003.07391}
\date{12 August 2026}

\begin{document}
\maketitle

\noindent\textbf{Status.}
Candidate counterexample, likely valid.  It gives a complete negative answer to
Section~7, item~(6) of the source.  The other open problems in that section are
not addressed.

\section{Source question}

Haddad, Jim\'enez, and Montenegro prove on a bounded open set
$\Omega\subset\R^n$ the reverse affine Zhang inequality
\[
 \Ep f\geq C_{n,p}(\Omega)
 \|f\|_p^{(n-1)/n}\|\nabla f\|_p^{1/n}.
\]
In Section~7, item~(6), PDF page~27 of \cite{source}, they ask whether this can
be improved by replacing $\|f\|_p$ with $\|f\|_q$ for some $q>p$.

\begin{center}
\includegraphics[width=\textwidth]{figures/open_problem_crop.png}
\end{center}

\section{Result}

\begin{theorem}\label{thm:main}
Let $n\geq2$, $p\geq1$, $q>p$, and let $\Omega\subset\R^n$ be any nonempty
open set.  Then
\[
 \inf_{0\neq f\in C_c^\infty(\Omega)}
 \frac{\Ep f}
 {\|f\|_q^{(n-1)/n}\|\nabla f\|_p^{1/n}}=0.
\]
Consequently there is no positive constant, even one depending on
$n,p,q,$ and $\Omega$, for the proposed $q>p$ inequality.
\end{theorem}

This shows that the exponent $q=p$ in the source theorem is sharp against
replacement by every strictly larger Lebesgue exponent.

\section{Proof intuition}

Affine energy is designed to tolerate volume-preserving anisotropy.  Compress a
fixed bump in only one coordinate.  Its ordinary gradient sees the inverse
thickness $\varepsilon^{-1}$, while its affine energy sees only the determinant
of the compression.  The $L^q$ norm contributes a different Jacobian power.
After inserting all three exact scalings, the proposed ratio contains the
positive power
\(
\varepsilon^{\frac{n-1}{n}(1/p-1/q)}
\).
It therefore vanishes for every $q>p$.  The construction is smooth and stays
inside the original fixed domain.

\section{Proof}

Choose $x_0\in\Omega$ and a rectangular neighborhood
$x_0+(-r,r)^n\subset\Omega$.  Fix a nonzero
$g\in C_c^\infty((-r,r)^n)$ such that $\partial_1g\not\equiv0$.  For
$0<\varepsilon\leq1$, let
\[
 A_\varepsilon=\operatorname{diag}(\varepsilon,1,\ldots,1),
 \qquad
 f_\varepsilon(x)=g\bigl(A_\varepsilon^{-1}(x-x_0)\bigr).
\]
Then $f_\varepsilon\in C_c^\infty(\Omega)$.

We first derive the affine-energy scaling.  For a smooth $f$, set
\[
 N_{f,p}(\xi)=\|\partial_\xi f\|_p,
 \qquad
 L_{f,p}=\{\xi\in\R^n:N_{f,p}(\xi)\leq1\}.
\]
The source's geometric formula is
\(
\Ep f=c_{n,p}n^{-1/n}|L_{f,p}|^{-1/n}
\).
For an invertible matrix $A$ and $f_A(x)=f(A^{-1}x)$, the chain rule and change
of variables give
\begin{align*}
 N_{f_A,p}(\xi)^p
 &=\int_{\R^n}
   |\langle A^{-\mathsf T}\nabla f(A^{-1}x),\xi\rangle|^p\,dx\\
 &=|\det A|\,N_{f,p}(A^{-1}\xi)^p.
\end{align*}
Therefore
\[
 L_{f_A,p}=A\bigl(|\det A|^{-1/p}L_{f,p}\bigr),
 \qquad
 |L_{f_A,p}|=|\det A|^{1-n/p}|L_{f,p}|,
\]
and hence the exact covariance law is
\begin{equation}\label{eq:affine-scale}
 \Ep(f_A)=|\det A|^{1/p-1/n}\Ep(f).
\end{equation}
Translation does not change any of the quantities.  Since
$\det A_\varepsilon=\varepsilon$, \eqref{eq:affine-scale} yields
\begin{equation}\label{eq:Escale}
 \Ep f_\varepsilon
 =\varepsilon^{1/p-1/n}\Ep g.
\end{equation}

The ordinary norms satisfy
\begin{equation}\label{eq:Lqscale}
 \|f_\varepsilon\|_q
 =\varepsilon^{1/q}\|g\|_q
\end{equation}
and, using only the first partial derivative,
\begin{equation}\label{eq:Gscale}
 \|\nabla f_\varepsilon\|_p
 \geq\|\partial_1f_\varepsilon\|_p
 =\varepsilon^{1/p-1}\|\partial_1g\|_p.
\end{equation}
Combining \eqref{eq:Escale}--\eqref{eq:Gscale} gives
\begin{align*}
 \frac{\Ep f_\varepsilon}
 {\|f_\varepsilon\|_q^{(n-1)/n}
  \|\nabla f_\varepsilon\|_p^{1/n}}
 &\leq C_g\,
 \varepsilon^{1/p-1/n-(n-1)/(nq)-(1/p-1)/n}\\
 &=C_g\,
 \varepsilon^{\frac{n-1}{n}(1/p-1/q)}.
\end{align*}
The last exponent is strictly positive because $q>p$.  Letting
$\varepsilon\downarrow0$ proves the theorem. \qed

\section{Checks, sharpness, and novelty status}

The argument uses identities rather than asymptotic equivalences.  It also
covers $p=1$: then the compressed first derivative has constant $L^1$ norm,
while the final exponent remains positive for every $q>1$.  Because the bump
is smooth, no approximation of characteristic functions or BV relaxation is
needed.

The accompanying script checks the exponent simplification exactly with
rational arithmetic and evaluates four representative parameter triples:
\begin{verbatim}
python code/verify_scaling.py
\end{verbatim}
All checks pass.  The script verifies algebra, not the analytic proof.

A bounded novelty search on 12 August 2026 used arXiv:2003.07391, the exact
title, and combinations of ``reverse Zhang'', ``affine Poincare'', ``q>p'',
and ``affine energy''.  It found the source and adjacent later affine Sobolev
work, but no explicit resolution of item~(6).  Novelty confidence is moderate,
pending a broader citation review.

\medskip
\noindent\textbf{Human-review recommendation.}
Verify the directional-unit-ball transformation and the determinant exponent
in \eqref{eq:affine-scale}.  Once that identity is confirmed, the remaining
norm scalings and the contradiction are immediate.

\begin{thebibliography}{9}
\bibitem{source}
J.~Haddad, C.~H. Jim\'enez, and M.~Montenegro,
\emph{From affine Poincar\'e inequalities to affine spectral inequalities},
arXiv:2003.07391 (2020), Section~7, item~(6), PDF p.~27.
\end{thebibliography}

\end{document}
